\documentclass[11pt,a4paper]{article}
\usepackage{templates/modeling-notes}
\hypersetup{pdftitle={Modeling patterns -- blending, networks, objectives and goals}}
\title{\textbf{Mathematical Modeling}\\[2mm]
\large Modeling Patterns: Blending, Networks, Objectives, and Goals}
\author{Personal study notes}
\date{}
\begin{document}
\maketitle
{\small\tableofcontents}

\section{Building an LP from a statement}\label{sec:principles}
A linear model expresses decisions using fixed coefficients, additive contributions,
and explicitly stated domains. A useful formulation requires more than recognizing
a familiar matrix: explain why a capacity is an inequality, a balance is an equality,
and an auxiliary variable preserves the intended criterion. The examples below
use continuous nonnegative quantities unless a different domain is stated.
All solutions are justified by hand bounds before numerical checks.

\section{Blending: a ratio that yields linear constraints}\label{sec:blending}
\subsection{Decisions, units, and conservation}
Choose monthly tonnes $x_1,\ldots,x_5\geq0$ of five input oils and tonnes
$y\geq0$ of finished oil. Assume no refining mass loss, negligible refining cost,
and a quality measure that is a mass-weighted average. The data are
\begin{center}
  \begin{tabular}{lrrrrr}\toprule
    Input & VEG1 & VEG2 & OIL1 & OIL2 & OIL3\\\midrule
    Purchase cost (PLN/t) & 110 & 120 & 130 & 110 & 115\\
    Hardness & $8.8$ & $6.1$ & $2$ & $4.2$ & $5$\\\bottomrule
  \end{tabular}
\end{center}
Finished oil sells for $150$ PLN/t. Vegetable refining capacity is $200$ t/month;
other-oil capacity is $250$ t/month. Conservation and capacities give
\[
  y=\sum_{i=1}^5x_i,\qquad x_1+x_2\leq200,\qquad x_3+x_4+x_5\leq250.
\]
The equality asserts that every purchased tonne enters the finished product;
it would require modification if some input were discarded or mass were lost.
Revenue minus purchase cost is
\[
  \max\ 150y-110x_1-120x_2-130x_3-110x_4-115x_5.
\]

\subsection{Why multiplying preserves the quality requirement}
For $y>0$, average hardness is
\[
  H=\frac{8.8x_1+6.1x_2+2x_3+4.2x_4+5x_5}{y}.
\]
Requiring $3\leq H\leq6$ is equivalent, for positive $y$, to
\begin{equation}\label{eq:quality}
  3y\leq8.8x_1+6.1x_2+2x_3+4.2x_4+5x_5\leq6y.
\end{equation}
The directions are preserved because the multiplier $y$ is positive. Every
coefficient is fixed, so these are linear inequalities. Multiplication by an
expression of unknown sign would need separate cases and would not justify
this step as written.

If $y=0$, conservation and $x\geq0$ force $x=0$. The average is undefined,
but the linear inequalities accept this no-production plan. Allowing no production
is a modeling convention. A delivery requirement $y\geq q>0$ would exclude it;
silently imposing $y>0$ instead changes the feasible set and its closedness.

As a simple illustration, mix $v$ tonnes of VEG1 and $o$ tonnes of OIL2.
The upper-quality bound becomes $2.8v\leq1.8o$, or $v/o\leq9/14$ when $o>0$.
An individual input above hardness $6$ can therefore be part of an acceptable mixture.
\begin{figure}[H]
  \centering
  \begin{tikzpicture}[x=1cm,y=1cm,font=\small]
    \draw[model axis] (0,0)--(9.3,0) node[right] {hardness};
    \draw[line width=3pt,modelblue!40] (3,0)--(6,0);
    \foreach \x/\lab in {3/3,4.2/4.2,6/6,8.8/8.8}
    \draw (\x,.12)--(\x,-.12) node[below] {$\lab$};
    \node[above] at (4.2,.15) {OIL2};
    \node[above] at (8.8,.15) {VEG1};
    \node[below,align=center] at (4.5,-.6) {Accepted mixture hardness\\lies in the blue interval.};
  \end{tikzpicture}
  \caption{Weighted averages interpolate between component qualities.
    At the upper boundary the VEG1 mass fraction is $9/23$.}\label{fig:hardness}
\end{figure}

\subsection{Eliminating a variable and proving the optimum}
Substitute $y=\sum_i x_i$ into profit. The margins become
$(40,30,20,40,35)$ PLN/t. In the upper-quality inequality, subtract $6y$:
\[
  \frac{14}{5}x_1+\frac1{10}x_2-4x_3-\frac95x_4-x_5\leq0.
\]
An optimal plan is
\begin{equation}\label{eq:blend-optimum}
  x=\left(\frac{4300}{27},\frac{1100}{27},0,250,0\right),\qquad y=450.
\end{equation}
Both refining capacities and the upper hardness bound are tight, giving $H=6$;
the lower bound is satisfied. Its profit is $475000/27$ PLN.

To prove optimality, weight vegetable capacity by $800/27$, other-oil capacity
by $140/3$, and the upper-quality inequality by $100/27$. All weights are
nonnegative. The combined left-side coefficients are
\[
  (40,30,860/27,40,1160/27),
\]
each at least the corresponding margin. Thus every feasible plan satisfies
\[
  40x_1+30x_2+20x_3+40x_4+35x_5
  \leq\frac{800}{27}\,200+\frac{140}{3}\,250=\frac{475000}{27}.
\]
Plan~\eqref{eq:blend-optimum} attains the bound. Strict excess coefficients for
$x_3,x_5$ force them to zero at every optimum; the three tight weighted rows
then determine the other amounts uniquely. The proof also explains why buying
the cheapest input independently of mixture quality can miss the optimum.

\section{Transportation with exact demand and unavailable routes}\label{sec:transport}
Let $x_{ij}\geq0$ be units sent from supplier $i$ to customer $j$. Use a variable
only for an allowed route $(i,j)\in E$. A prohibited route is absent, or its flow
is fixed to zero; an infinite cost is a shorthand for prohibition, not a finite
coefficient to pass to an LP solver.
\[
  \min\sum_{(i,j)\in E}c_{ij}x_{ij},\qquad
  \sum_{j:(i,j)\in E}x_{ij}\leq S_i,\qquad
  \sum_{i:(i,j)\in E}x_{ij}=D_j.
\]
Supplier capacities are upper bounds. Customer demand is exact in this model;
if oversupply were allowed, the appropriate customer condition could instead
be a lower bound. Total capacity at least total demand is necessary, but route
restrictions can make particular customer subsets impossible to serve.

\subsection{A complete small instance and a lower-bound certificate}
Take $S=(135,56,93)$ and $D=(62,83,39,91)$. Costs per unit are
\begin{center}
  \begin{tabular}{lrrrr}\toprule
    & Customer 1 & Customer 2 & Customer 3 & Customer 4\\\midrule
    Supplier 1 & 132 & --- & 97 & 103\\
    Supplier 2 & 85 & 91 & --- & ---\\
    Supplier 3 & 106 & 89 & 100 & 98\\\bottomrule
  \end{tabular}
\end{center}
The following flow matrix satisfies demands and capacities, with dashes denoting
unavailable routes:
\[
  X=\begin{pmatrix}0&\text{---}&39&87\\56&0&\text{---}&\text{---}\\6&83&0&4\end{pmatrix}.
\]
Row totals are $(126,56,93)$; nine units of supplier 1 capacity are unused.
The cost is $25919$.

Choose customer values $v=(111,94,97,103)$ and supplier charges $u=(0,26,5)$.
For each allowed route, check $c_{ij}\geq v_j-u_i$. Multiplying by its
nonnegative flow and summing gives
\[
  \sum c_{ij}x_{ij}\geq\sum_j v_jD_j-\sum_i u_i\sum_jx_{ij}
  \geq\sum_jv_jD_j-\sum_i u_iS_i=25919.
\]
The second inequality uses $u\geq0$ and the supplier upper bounds. The proposed
plan reaches this lower bound. Used routes have equality in the route bound,
and each supplier with a positive charge uses its full capacity; these facts
explain the certificate's equality conditions.

\section{Production with storage as a forward-time network}\label{sec:storage}
\subsection{Why the transportation costs contain time}
Let $p_{k,d}$ and $p_{k,n}$ be day- and night-shift production in month $k$.
Let $I_k\geq0$ be stock after meeting month $k$'s demand $D_k$. Assume no
backlogging, no mass loss, unlimited storage, initially empty stock, and no
desired final surplus: $I_0=I_4=0$. Conservation gives
\begin{equation}\label{eq:inventory}
  I_{k-1}+p_{k,d}+p_{k,n}-D_k=I_k.
\end{equation}
The data are
\begin{center}
  \begin{tabular}{lrrrr}\toprule
    Month & January & February & March & April\\\midrule
    Day capacity & 100 & 150 & 140 & 160\\
    Night capacity & 50 & 75 & 70 & 80\\
    Demand & 80 & 200 & 300 & 200\\\bottomrule
  \end{tabular}
\end{center}
Day production costs $1$ PLN/unit, night production $1.5$, and each unit stored
for one month costs $0.30$. The objective is
\[
  \min\sum_{k=1}^4(p_{k,d}+1.5p_{k,n})+0.30\sum_{k=1}^3 I_k,
\]
with production between zero and its capacity.

Equivalently, let $x_{k,s,t}$ be production from shift $s$ in month $k$ used
for demand in month $t$. Only $t\geq k$ is allowed, since later production
cannot meet earlier demand. If shift $s$ costs $a_s$, that route costs
$a_s+0.30(t-k)$ per unit. Supplier rows limit production and customer rows
equal demand. Each unit assigned to a later month crosses $t-k$ storage
intervals, proving the equivalence of route costs and inventory holding cost.
\begin{figure}[H]
  \centering
  \begin{tikzpicture}[font=\small]
    \foreach \k/\x in {1/0,2/3,3/6,4/9} {
    \node[draw,rounded corners] (m\k) at (\x,0) {Month $\k$};
    \draw[model flow] (\x,1.1)--(m\k) node[midway,right] {$p_{\k,d}+p_{\k,n}$};
    \draw[model flow] (m\k)--(\x,-1.1) node[midway,right] {$D_{\k}$};
    }
    \draw[model flow] (m1)--(m2) node[midway,above] {$I_1$};
    \draw[model flow] (m2)--(m3) node[midway,above] {$I_2$};
    \draw[model flow] (m3)--(m4) node[midway,above] {$I_3$};
  \end{tikzpicture}
  \caption{Stock carries material only forward. Each node's inflow and production
    equal its demand and outgoing stock.}\label{fig:inventory}
\end{figure}

\subsection{An optimal schedule without a black-box solve}
An optimal schedule is
\begin{center}
  \begin{tabular}{lrrrr}\toprule
    Month & January & February & March & April\\\midrule
    Day production & 100 & 150 & 140 & 160\\
    Night production & 45 & 75 & 70 & 40\\
    Closing inventory & 65 & 90 & 0 & 0\\\bottomrule
  \end{tabular}
\end{center}
Check each balance in~\eqref{eq:inventory}. Total demand is $780$, while total
day capacity is $550$. Any admissible schedule therefore needs at least $230$
night units and production cost at least $780+0.5\cdot230=895$ PLN.

March demand is $300$, but March capacity is $210$, so $I_2\geq90$.
Combining February and March balances with their capacities gives
\[
  I_1\geq200+300-(225+210)=65.
\]
Other inventories are nonnegative. Holding cost is therefore at least
$0.30(65+90)=46.5$ PLN. The displayed schedule attains both bounds, giving
the exact optimum $941.5=1883/2$ PLN. Enough total capacity alone would not
prove feasibility: these prefix-demand arguments also respect time.

\section{Minimum-cost flow and the meaning of signs}\label{sec:flow}
For a directed graph $(V,E)$ choose flow $f_{uv}\geq0$ on each arc, bounded by
capacity $C_{uv}$. Let $s_v$ be net supply: positive at a source and negative
at a demand node. The model is
\[
  \min\sum_{(u,v)\in E}p_{uv}f_{uv},\qquad 0\leq f_{uv}\leq C_{uv},\qquad
  \sum_{u:(u,v)\in E}f_{uv}+s_v=\sum_{w:(v,w)\in E}f_{vw}.
\]
Incoming flow plus local production equals outgoing flow. Summing all node
balances cancels every arc and gives $\sum_v s_v=0$, a necessary condition.
It is insufficient when capacities or missing routes block deliveries.
With an incidence matrix having $+1$ at an arc's tail and $-1$ at its head,
the same convention is $Bf=s$. Reversing the matrix convention also reverses
the sign of the supply vector; keeping only one of these changes is an error.

\begin{figure}[H]
  \centering
  \begin{tikzpicture}[font=\small]
    \node[circle,draw] (a) at (0,0) {$A$};
    \node[circle,draw] (b) at (3,1.5) {$B$};
    \node[circle,draw] (c) at (6,0) {$C$};
    \draw[model flow] (a)--(b) node[midway,above left] {cost 1; cap 2};
    \draw[model flow] (b)--(c) node[midway,above right] {cost 1; cap 2};
    \draw[model flow] (a)--(c) node[midway,below] {cost 5; cap 3};
    \node[below] at (0,-.5) {$s_A=3$};
    \node[above] at (3,2) {$s_B=0$};
    \node[below] at (6,-.5) {$s_C=-3$};
  \end{tikzpicture}
  \caption{Three units must reach $C$. The cheap path has a bottleneck of two units.}
  \label{fig:flow}
\end{figure}
Here balances force $f_{AB}=f_{BC}=t$ and $f_{AC}=3-t$, with $0\leq t\leq2$.
Cost is $2t+5(3-t)=15-3t$, minimized at $t=2$ with value $9$.
This derivation simultaneously checks conservation, capacities and optimality.

\section{Choosing an objective and comparing several goals}\label{sec:objectives}
\subsection{Constants, feasibility, and genuine decisions}
An unavoidable fixed cost $K$ changes profit $r^Tx$ to $r^Tx-K$. It changes
the reported value but leaves the optimizers unchanged. A startup cost incurred
only if production occurs is decision-dependent: modeling it normally introduces
an activation decision, often a binary variable. It cannot be dropped as an
unavoidable constant.

If the question is simply whether requirements can be met, use an objective
such as $\min0$ over the original feasible set. Every feasible point then has
equal objective value. Also, $x_1+x_2\leq2$, $x_1\geq1$, $x_2\geq1$ force
$(x_1,x_2)=(1,1)$ independently of the objective. Recognizing this singleton
does not justify dropping sign or inequality constraints in a general model.

\subsection{Weights, thresholds, and max--min reformulation}
Suppose several affine criteria $f_i(x)$ are to be maximized. Their individual
optima help reveal trade-offs, but may be incompatible. Three LP formulations are:
\begin{enumerate}[leftmargin=1.8em]
  \item \textbf{Weighted sum:} maximize $\sum_i w_i f_i(x)$ with nonnegative
        weights. Scaling and units are part of the preference specification.
  \item \textbf{Thresholds:} maximize one criterion subject to $f_i(x)\geq\beta_i$
        for the others. Thresholds may make the model infeasible.
  \item \textbf{Worst normalized achievement:} maximize $\min_i g_i(x)$, where
        the $g_i$ are comparable affine normalized criteria. Introduce a free
        variable $z$ and maximize $z$ subject to $z\leq g_i(x)$ for every $i$.
\end{enumerate}
For fixed $x$, the largest feasible $z$ is exactly $\min_i g_i(x)$, proving
the last reformulation. Imposing $z\geq0$ without justification could exclude
valid decisions with negative scores. Positive weights ensure that an attained
weighted-sum optimum is Pareto efficient: a feasible decision improving one
criterion and worsening none would strictly improve the weighted sum.

For $0\leq x\leq1$, take $f_1=x$ and $f_2=1-x$. Equal weights make every
point optimal. A greater weight on $f_1$ selects $x=1$; a greater weight on
$f_2$ selects $x=0$. A threshold $f_2\geq\beta$, $0\leq\beta\leq1$, gives
$x^*=1-\beta$ when maximizing $f_1$. The max--min model has
$z\leq x$, $z\leq1-x$, hence $2z\leq1$; $x=z=1/2$ attains this bound.
Replacing $f_1$ by $100x$ without adjusting normalization changes these choices.

\section{Soft constraints: solving a deliberately changed model}\label{sec:goals}
\subsection{What the deviation variables measure}
Keep inviolable conditions in a hard feasible set. For a relaxable upper target
$a^Tx\leq b$, use $a^Tx\leq b+u$, $u\geq0$. For a lower target use
$a^Tx\geq b-v$, $v\geq0$. For an equality target use
$a^Tx=b+u-v$ with $u,v\geq0$. With positive penalties, minimizing these
variables gives one-sided violations or, for the equality, $|a^Tx-b|$.
With a strictly positive weighted-sum penalty, the equality deviations cannot
both be positive at an optimum: reducing both
by their minimum preserves the equality and reduces the sum penalty.

Use either a weighted sum of deviations or a variable $z\geq0$ bounding
all weighted deviations to minimize the largest one. Deviations in tonnes,
hours and PLN are incomparable without units-aware weights or normalization.
Conservation laws and domains should stay hard unless changing them has a
deliberate meaning. If the hard set itself is empty, softening other targets
cannot restore feasibility.

\subsection{A hand comparison of two relaxation criteria}
Let $x\geq0$ be hard and request $x\leq1$, $x\geq3$. These targets are
incompatible. Relax them to
\[
  x\leq1+u,\qquad x\geq3-v,\qquad u,v\geq0.
\]
Subtracting the inequalities gives $u+v\geq2$. Minimizing $u+v$ has value
$2$, attained by every $1\leq x\leq3$ with $u=x-1$, $v=3-x$.
Minimizing $z$ with $u,v\leq z$ gives $2z\geq u+v\geq2$; the unique
balanced answer is $x=2$, $u=v=z=1$. Weighting violations differently,
such as minimizing $2u+v$, selects $x=1$, $u=0$, $v=2$ with value $2$.
Thus the choice of violation measure determines what compromise means.

Phase I in simplex also uses nonnegative artificials, but its purpose is to
test feasibility of the original hard equations. A positive \emph{minimum}
artificial sum, equivalently a negative maximum of minus that sum, certifies
failure of that test. Goal programming instead accepts
specified violations and reports a solution of a modified problem.

\section{Questions for checking understanding}\label{sec:review}
\begin{enumerate}[leftmargin=1.8em]
  \item Why does multiplying a quality ratio by output preserve inequality directions?
  \item What does the blending model mean at zero production?
  \item Why should unavailable routes be omitted rather than assigned a huge cost?
  \item Why is total production capacity insufficient to guarantee on-time delivery?
  \item How do node-balance signs depend on the incidence-matrix convention?
  \item What changes when an objective is multiplied by a conversion factor?
  \item Which requirements can be relaxed without losing the physical meaning of a model?
\end{enumerate}

\section{Sources and editorial remarks}\label{sec:sources}
The supplied current-course source is \texttt{ProgMat\_w1.pdf}, 39 PDF pages,
stored externally in \texttt{../lectures-2026-27/}. Its delivery date is not
established. This document develops the modeling material in pp.~10--26:
blending in Section~\ref{sec:blending}, transportation in Section~\ref{sec:transport},
storage in Section~\ref{sec:storage}, flow in Section~\ref{sec:flow}, objective
selection in Section~\ref{sec:objectives}, and goals in Section~\ref{sec:goals}.
The supplied numerical data are retained. Hand optima, weighted bounds, the
three-node example, normalized-goal example, relaxation example, and diagrams
are added explanations. Continuous quantities, empty initial/final stock, no
backlogging and unlimited storage are explicit modeling assumptions.

This conceptual supplement keeps the earlier LP geometry and proofs in
Chapter~2 intact. The course map records reading routes
and current task requirements separately from historical enrichment.

\paragraph{Numerical checks.}
The actual companion is
\href{run:matlab/exercises/modeling_patterns_checks.m}{modeling-patterns MATLAB script}
at \path{matlab/exercises/modeling_patterns_checks.m}.
It checks the same factory, blending, shipment, storage, flow, regression, diet,
objective and relaxation examples using base MATLAB and absolute tolerance
$10^{-10}$. Expected values include blending $475000/27$, transportation $25919$,
storage $1883/2$, and flow $9$. Exact rational calculations independently
checked the new hand results. MATLAB execution of the new companion is unverified;
the execution record is \path{docs/matlab-verification.md}.
\end{document}
